% ANCHOR: after_label def:c9-bases
\begin{remark}[Physics notation: the three-leg functionals in components]
In the components of the physics notation key after Section 1.2, with $F^0_{mn}=\partial_mA_n-\partial_nA_m$, $\partial^2=\partial_m\partial_m$, $A_{ij}=A_I\Sigma^I_{ij}$, $A^{ij}=\tfrac12\varepsilon^{ijkl}A_{kl}$ and $(-,-)$ on the colour factors,
\[
\begin{aligned}
&G_1=-\int d^4x\,\big([c,A_n],\partial_mF^0_{mn}\big),\qquad G_2=-\int d^4x\,\big([c,A_n],\partial_n\partial_mA_m\big)=-\int d^4x\,\big(\partial_mA_m,\partial_n[A_n,c]\big),\\
&G_3=-\sum_I\int d^4x\,\big([c,A_I],\partial^2A_I\big)=-\tfrac14\int d^4x\,\big([c,A^{ij}],\partial^2A_{ij}\big),\qquad G_4=\int d^4x\,\big([c,\partial_mc],A^*_m\big),\\
&H_1=2\int d^4x\,\big(\partial_mA_n,[A_m,A_n]\big),\qquad H_2=\sum_I\int d^4x\,\big(\partial_mA_I,[A_m,A_I]\big)=\tfrac14\int d^4x\,\big(\partial_mA^{ij},[A_m,A_{ij}]\big),\\
&H_3=\int d^4x\,\big([A_m,c],A^*_m\big),\qquad H_4=\tfrac14\int d^4x\,\big([A^{ij},c],A^*_{ij}\big),\qquad H_6=\int d^4x\,\big([c,c],c^*\big),
\end{aligned}
\]
and $G_5$, $H_5$ and the two Yukawa-type functionals are written in two-component form in Remark \ref{rem:c9-twocomp3}. The $G_i$ have ghost number $1$: $G_1$ and $G_2$ are the transverse and the longitudinal ghost--gluon--gluon structures, $G_3$ the ghost--scalar--scalar, $G_5$ the ghost--gaugino--gaugino and $G_4$ the ghost--ghost--gluon-antifield structure. The $H_j$ have ghost number $0$: $H_1$ and $H_2$ are the three-gluon and gluon--scalar--scalar vertices, and $H_3,\dots,H_6$ the antifield (BRST-source) vertices, the cubic tree vertices written out at the end of Section 4.3.1.
\end{remark}

% ANCHOR: after_label prop:c9-basis
\begin{remark}[Physics notation: leg sets, colour and parity]
The four leg sets of Lemma \ref{lem:c9-legs} are ghost--gluon--gluon $(c,A_m,A_n)$, ghost--scalar--scalar $(c,A_I,A_J)$, ghost--gaugino--gaugino $(c,\lambda,\lambda)$ and two ghosts with the gluon antifield $(c,c,A^*_m)$: the three-leg local functionals of ghost number $1$, where the cubic part of the breaking of the Slavnov--Taylor identity lies. For $\mathfrak{su}(N)$, with each leg written $X_j=-igX_j^aT^a$ as in the key, the colour factor is
\[
f=([X_0,X_1],X_2)=2\hbar g\,f^{abc}X_0^aX_1^bX_2^c,\qquad \operatorname{tr}_{\mathfrak g}(\mathrm{ad}_X\mathrm{ad}_Y\mathrm{ad}_Z)=-\tfrac{g^3N}2\,f^{abc}X^aY^bZ^c,
\]
the second by $\operatorname{tr}(T^a_AT^b_AT^c_A)=\tfrac i2N f^{abc}$ for the adjoint generators $(T^a_A)^{bc}=-if^{abc}$. The plane-wave values of Proposition \ref{prop:c9-basis} are those of the components after Definition \ref{def:c9-bases}, with incoming momenta and $\partial_m\mapsto ip_m$. The parity-odd structure of $(c,\lambda,\lambda)$, written in two-component form in Remark \ref{rem:c9-twocomp3}, is the $\varepsilon_{mnkl}$-type ($\gamma_5$-type) structure; its coefficient vanishes because the reflection of one spacetime and one transverse direction is a symmetry of the regularised theory.
\end{remark}

% ANCHOR: after_label cor:c9-completions
\begin{remark}[Physics notation: linearised BRST images, cubic couplings and completions]
$Q$ acts on the $H_j$ as the linearised BRST transformation with the arrows at the end of Section 1.5, and for a functional $X$ of $A$ alone $\{I,X\}$ is minus the change of $X$ under the non-abelian part $A_m\mapsto[A_m,c]$, $A_I\mapsto[A_I,c]$ of the gauge transformation; thus $\{I,\int(\partial_mA_m,\partial_nA_n)\}=2G_2$ comes from the Feynman-gauge term. The four kernel directions of Proposition \ref{prop:c9-maps} (ii) are, with the cubic part of $I$ of that proposition,
\[
\begin{aligned}
&H_1-4H_3-2H_6=-4\big(-\tfrac14H_1+H_3+\tfrac12H_6\big),\qquad H_2-2H_4=-2\big(-\tfrac12H_2+H_4\big),\\
&\int(\lambda,\rho([A|_V,\lambda]))+2H_5=2\Big(\tfrac12\int(\lambda,\rho([A|_V,\lambda]))+H_5\Big),\qquad \int(\lambda,\rho([A|_{W_{\mathbb R}^\perp},\lambda])):
\end{aligned}
\]
the three-gluon coupling with its ghost and antifield vertices, the gluon--scalar coupling with its antifield vertex, the gluon--gaugino coupling with its antifield vertex, and the Yukawa coupling, each invariant under $Q$. The completions of Corollary \ref{cor:c9-completions} are the gauge-covariant kinetic terms $\int(F_{A|_V},F_{A|_V})=\tfrac12\int d^4x\,(F_{mn},F_{mn})$, $\int(d_{A|_V}A|_{W_{\mathbb R}^\perp},d_{A|_V}A|_{W_{\mathbb R}^\perp})=\tfrac14\int d^4x\,(D_mA^{ij},D_mA_{ij})$, twice the gaugino terms of $S$ and the antifield part of $S$, written out after Definition \ref{def:c1-action}.
\end{remark}

% ANCHOR: after_label thm:c9-ob3
\begin{remark}[Physics notation: the three-point breaking of the Slavnov--Taylor identity]
\eqref{eq:c9-ob3} is the one-loop breaking of the Slavnov--Taylor identity by the regulator in the three-point sectors, in the components after Definition \ref{def:c9-bases}; for $\mathfrak{su}(N)$ the unit is $-\frac{g^2N}{32\pi^2}$ and each $G_i$ carries the colour factor $2\hbar g\,f^{abc}X_0^aX_1^bX_2^c$, so
\[
\hbar\,\mathrm{Ob}^{(0)}=-\frac{g^2N}{32\pi^2}\Big[\big(\tfrac{11}4+\tfrac{19}4\log\tfrac32\big)(G_1+G_3)+\tfrac{27}8G_2-9\log\tfrac32\,G_5\Big]
\]
in the three-leg sectors: the transverse ghost--gluon--gluon and the ghost--scalar--scalar structures with one coefficient, the longitudinal one with $\tfrac{27}8$, the ghost--gaugino--gaugino structure with $-9\log\tfrac32$, and no breaking with two ghosts and the gluon antifield. The cycle graphs of the table are one-loop triangles and bubbles with one special edge $K_\varepsilon-K_0$, the smeared minus the sharp field--antifield contraction (end of Section 4.4); the bridge graphs join a cubic vertex to a one-loop self-energy subgraph by the special edge, and the reducible graphs are the two-point breaking of Chapter 4 attached to a cubic vertex by a propagator.
\end{remark}

% ANCHOR: after_label thm:c9-F3
\begin{remark}[Physics notation: the cubic finite counterterm]
$\hbar F_1^{(0)}$ is a finite symmetry-restoring counterterm in the sense of algebraic renormalisation, entering the interaction as $I-\hbar F_1^{(0)}$. The five equations are those of the coefficients of $G_1,\dots,G_5$; $G_2$ is reached by the Feynman-gauge term alone. The free directions $\beta_k$ are finite multiples of the gauge-covariant kinetic terms and of the antifield part of $S$ (up to the $\kappa_k$ directions), and the $\kappa_k$ finite changes of the four cubic couplings (block after Corollary \ref{cor:c9-completions}). For $\mathfrak{su}(N)$, with
\[
\int(dA,dA)=-4\hbar\int d^4x\,\Big(-\tfrac14F^{0a}_{mn}F^{0a}_{mn}-\tfrac12\partial_mA^a_I\partial_mA^a_I\Big),\qquad \int(\lambda,{\not d}\lambda)=2\hbar\int d^4x\,i\,\bar\lambda^a_{\dot\alpha i}\bar\sigma^{m\dot\alpha\alpha}\partial_m\lambda^a_\alpha{}^i,
\]
and $\int(d^*A|_V,d^*A|_V)=2\hbar\int d^4x\,(\partial_mA^a_m)^2$, the representative \eqref{eq:c9-F2rep}, with vanishing cubic part, is
\[
\frac{\hbar F_1^{(0)}}\hbar=\int d^4x\Big\{-\frac{27g^2N}{256\pi^2}(\partial_mA^a_m)^2+\frac{g^2N}{8\pi^2}\Big(\tfrac{11}8+\tfrac{19}8\log\tfrac32\Big)\Big(-\tfrac14F^{0a}_{mn}F^{0a}_{mn}-\tfrac12\partial_mA^a_I\partial_mA^a_I\Big)+\frac{9g^2N}{32\pi^2}\log\tfrac32\;i\,\bar\lambda^a\bar\sigma^m\partial_m\lambda^a\Big\}
\]
up to terms with four or more legs: a finite change of the Feynman-gauge term, of the free gluon and scalar Lagrangian and of the free gaugino Lagrangian.
\end{remark}

% ANCHOR: after_label thm:c9-wz
\begin{remark}[Physics notation: the Wess--Zumino consistency condition]
$s\,\mathrm{Ob}=0$ is the Wess--Zumino consistency condition of the one-loop anomaly (block after Theorem \ref{thm:c8-obstruction}). In the three-leg sector it ties the longitudinal ghost--gluon--gluon coefficient $\alpha_2$ to the two-point breaking $-\tfrac{27}8$ units times $\int d^4x\,(\partial^2c,\partial_mA_m)$ of Chapter 4, through
\[
QG_2=-\int d^4x\,\big([c,\partial_nc],\partial_n\partial_mA_m\big)-\int d^4x\,\big([c,A_n],\partial_n\partial^2c\big)
\]
on the sector $(c,c,A_m)$: the condition reads $(\alpha_2-\tfrac{27}8)\,QG_2=0$, it leaves the transverse, scalar, gaugino and $(c,c,A^*_m)$ coefficients free, and it is the condition for the three-point breaking to be removable by the finite counterterm of Theorem \ref{thm:c9-F3}.
\end{remark}

% ANCHOR: end_of_section 9.2
\begin{remark}[Physics notation: the dilatation vertex and the composite $A_m$--$c$ line]
$S^{(1)}(\Delta)$ is the dilatation vertex written out at the end of Section 6.6, and $\{S^{(1)}(\Delta),F\}=v.F$ applies the classical scale transformation $2x\cdot\partial+2w$ to the fields; on a local operator of canonical dimension $w_O$, $v.O=2(x^m\partial_m+w_O)$ on its density. In \eqref{eq:c9-D1a}, $P\star v.$ attaches the vertex to one propagator, and the composite kernel $2(P\star v.P\star\otimes1)K_0$ is one line made of the gluon propagator, the dilatation vertex $(v.A,A^*)$ and the ghost--antifield propagator: it joins a leg $c$ to a leg $A_m$. It exists because $(Q^{GF})^2\neq0$ in Feynman gauge (end of Section 3.1): $(P\star)^2=(Q^{GF})^2\big(\int_\varepsilon^Te^{-t|p|^2}dt\big)^2$ maps $A^*_m$ to $-2ip_mA^*_m$ in the slot of $c$ and $c^*$ to $-2ip_mc^*$ in the slot of $A_m$, times $e^{-2\varepsilon|p|^2}/|p|^4$ at $T=\infty$; the line is longitudinal and of order $|p|^{-3}$ at small momentum. Proposition \ref{prop:c9-ward} states that the dilatation moves both endpoints of the proper-time integral, the ultraviolet cutoff $\varepsilon$ and the endpoint $T$ of the transfer of [C07], a Wilsonian scale; on the scalar factor of the $A\otimes A$ component, $\big(2x\cdot\partial_x+2y\cdot\partial_y+4\big)\frac{e^{-|x-y|^2/4T}-e^{-|x-y|^2/4\varepsilon}}{4\pi^2|x-y|^2}=-4\big(\varepsilon\partial_\varepsilon+T\partial_T\big)\frac{e^{-|x-y|^2/4T}-e^{-|x-y|^2/4\varepsilon}}{4\pi^2|x-y|^2}$. Corollary \ref{cor:c9-CS} is the Callan--Symanzik equation of $O[T]$ at the Wilsonian scale $T$, twice that of Theorem \ref{thm:c6-CS} at $T=\infty$ because $\Delta=2E$.
\end{remark}

% ANCHOR: after_label thm:c9-quadratic
\begin{remark}[Physics notation: the dilatation obstruction and its counterterm]
$\hbar\,\mathrm{Ob}^{(1)}(\Delta)$ is the one-loop breaking of the dilatation Ward identity by the regulator, and $\hbar F_1^{(1)}(\Delta)$ the finite counterterm of the dilatation vertex that restores it in the quadratic sector. In components, with $\int(dA,dA)$ and $\int(dc,A^*)=\int d^4x\,(\partial_mc,A^*_m)$ as after Theorem \ref{thm:c4-quadratic}, the field--antifield pairings after Definition \ref{def:c1-fields} and the gaugino functionals of Remark \ref{rem:c9-twocompD},
\begin{itemize}
\item the arity-one counterterm (i) is $\log\varepsilon$ units times the image under $Q$ of $\int d^4x\,(x_mA_m,c^*)-\tfrac12\int(\rho(x)\lambda^*,\lambda^*)$, the counterterm of the composite cycles;
\item \eqref{eq:c9-ObDelta} is $-4$ units times the gluon and scalar kinetic functional $\int(dA,dA)$, $16$ units times $\int(\lambda,{\not d}\lambda)$ and $-4$ units times $\int(\partial_mc,A^*_m)$, the scaling of the regulator (Chapter 6), plus the composite part $\tfrac13$, $-\tfrac13$, $-\tfrac23$ units from the composite line;
\item for the translation-invariant solution (iii), the bracket with $S^{(1)}(\Delta)-\hbar F_1^{(1)}(\Delta)$ acts on the coordinates of $A_m$, $A_{ij}$, of $\lambda_\alpha{}^i$, $\bar\lambda_{\dot\alpha i}$ and of $c$ as $v.$ minus $\tfrac{22}3$, $\tfrac{47}3$ and $\tfrac83$ units times the field (Lemma \ref{lem:c6-count}): $4$ units times the field-counting weights $2$, $4$, $1$, minus $\tfrac23$, $\tfrac13$, $\tfrac43$ units from the composite line.
\end{itemize}
In \eqref{eq:c9-Fperp}, $\int\big(2E(A|_{W_{\mathbb R}^\perp})+2A|_{W_{\mathbb R}^\perp},A^*|_{W_{\mathbb R}^\perp}\big)=\sum_I\int d^4x\,\big((2x^m\partial_m+2)A_I,A^*_I\big)$, so $\theta$ multiplies the classical dilatation of the six scalars.
\end{remark}

% ANCHOR: after_label thm:c9-D1
\begin{remark}[Physics notation: the one-loop dilatation of the scalar bilinears]
For the physics dilatation $D=-E=-\tfrac12\Delta$ (the generator $-E$ of [SW, (3)]), acting through $-\Lambda^{(1)}$, the one-loop dilatation operator is $\tfrac12D_1$ (block after Definition \ref{def:c8-components}), and $\tfrac12\{S^{(1)}(\Delta),O\}=(x^m\partial_m+2)O$ is the classical scale transformation of a dimension-$2$ operator. On the part with two external scalars and for $\mathfrak{su}(N)$, \eqref{eq:c9-D1} reads
\[
\tfrac12D_1(O)=\Big(\tfrac13+\theta\Big)\frac{g^2N}{32\pi^2}\,(x^m\partial_m+2)O\quad(\mathbf{20'}),\qquad \tfrac12D_1(O)=\frac{3g^2N}{4\pi^2}\,O+\Big(\tfrac13+\theta\Big)\frac{g^2N}{32\pi^2}\,(x^m\partial_m+2)O\quad(\text{Konishi}),
\]
for the operators written out after Definition \ref{def:c5-bilinear}: the anomalous dimensions $0$ and $\tfrac{3g^2N}{4\pi^2}$ of Theorem \ref{thm:c6-dil} plus a one-loop multiple of the classical dilatation that the two-point equations leave free. The rows of the table are the regulated graphs of the dilatation Ward identity through $O$ and $S^{(1)}(\Delta)$ with one special edge, the finite counterterm along $\Delta$ on the two scalar legs, and the vanishing arity-one counterterm of $O$. The composite triangle of Proposition \ref{prop:c9-triangle} has on its cycle $O$, the gluon--scalar--scalar vertex, the vertex $([A,c],A^*)$ and the composite $A_m$--$c$ line (end of Section 9.2); its image under $Q$ is $-2$ units times the classical transformation.
\end{remark}

% ANCHOR: end_of_section 9.4
\begin{remark}[Physics notation: $\theta$, Konishi and the $\mathbf{20'}$]
Theorem \ref{thm:c9-kappa} fixes $\theta=-\tfrac13$ by the one-loop form of the commutator $[D,P_m]=P_m$, i.e. $[\Delta,P_b]=-2P_b$ (Remark \ref{rem:c2-dilatation}), on the scalar bilinears: a multiple $a$ of the classical dilatation in $D_1$ appears as $2a\,b^m\partial_mO$ in this commutator, while the right side, $s$ applied to the arity-two output, is linear in the position and has no such term. With $\theta=-\tfrac13$, on the part with two external scalars,
\[
[D,O(x)]=(x^m\partial_m+2+\gamma_O)\,O(x),\qquad \gamma_{\mathbf{20'}}=0,\qquad \gamma_{\rm Konishi}=\frac{3g^2N}{4\pi^2},
\]
for the twenty traceless operators $\operatorname{tr}(\phi_I\phi_J)-\tfrac16\delta_{IJ}\operatorname{tr}(\phi_K\phi_K)$ and the Konishi operator $\operatorname{tr}(\phi_I\phi_I)$, $\phi_I=A^a_IT^a$: the values of Theorem \ref{thm:c6-dil} as components of the one-loop equivariant structure, up to the parts of $D_1(O)$ with legs $(A_m,A_m)$ and $(c,A^*_m)$, which (iii) requires to be BRST-exact (in the image of $\Lambda^{(0)}$). The finite counterterm along $\Delta$ in the scalar sector is then $\big[\tfrac{20}3\sum_I\int d^4x\,(A_I,A^*_I)-\tfrac23\sum_I\int d^4x\,(x^m\partial_mA_I,A^*_I)\big]$ units. For every $\theta$, the coefficients of $O$ in $\tfrac12D_1(O)$ for the Konishi operator and for the $\mathbf{20'}$ differ by $\tfrac{3g^2N}{4\pi^2}$.
\end{remark}

% ANCHOR: end_of_section 9.5
\begin{remark}[Physics notation: the translation component and a finite renormalisation]
For $\mathfrak{su}(N)$, on the part with two external scalars, \eqref{eq:c9-P1} reads $P_1(O)=\big(\tfrac13+\theta_P\big)\frac{g^2N}{32\pi^2}\,b^m\partial_mO$, a one-loop multiple of the classical translation: the term $\tfrac13$ comes from the cycles in which $S^{(1)}(P_b)$ appears as the composite $A_m$--$c$ line $2\big(\partial_b(P\star)^2\otimes1\big)K_0$, with the derivative $\partial_b$ in place of the dilatation vertex, and $\theta_P$ from the finite counterterm $\theta_P$ units times $\sum_I\int d^4x\,(b^m\partial_mA_I,A^*_I)$. In Proposition \ref{prop:c9-scheme}, $\hbar\psi'=\tfrac a2\{S^{(1)}(\Delta),-\}$ sends an operator $O$ of canonical dimension $w_O$ to $a(x^m\partial_m+w_O)O$, $a$ times the infinitesimal Weyl scaling of Lemma \ref{lem:c6-data}: a finite renormalisation of the composite operators that changes the component along a generator of scaling dimension $d_x$ ($1$ for $P$, $\tfrac12$ for $Q$, $\bar Q$, $-\tfrac12$ for $S$, $\bar S$, $-1$ for $K$, $0$ for $D$, $M$, $\bar M$, $R$; end of Section 7.4) by $-a\,d_x$ times its classical transformation. With $a=-(\tfrac13+\theta_P)$ units it sets $P_1=0$ and leaves $D_1$ and the anomalous dimensions unchanged. The components of this chapter are computed at $T=\infty$ and are the same at every finite endpoint $T$ of the transfer (Remark \ref{rem:c9-endpoint}, Theorem \ref{thm:c8-L} (ii)); the composite line is of order $|p|^{-3}$ at small momentum, $|x-y|^{-1}$ in position space.
\end{remark}
